% Part: model-theory
% Chapter: models-of-arithmetic
% Section: computable-models

\documentclass[../../../include/open-logic-section]{subfiles}

\begin{document}

\olfileid{mod}{mar}{cmp}
\section{مدل‌های محاسبه‌پذیر حساب}

\begin{explain}
مدل استاندارد~$\Struct{N}$ دو ویژگی خوب دارد. دامنه‌اش اعداد
طبیعی~$\Nat$ است؛ یعنی عناصرش دقیقاً از همان نوع چیزهایی‌اند که
می‌خواهیم با زبان حساب درباره‌شان سخن بگوییم، و عددنمای استاندارد
$\num{n}$ واقعاً خود~$n$ را مشخص می‌کند. ویژگی خوب دیگر این است که
همهٔ تعبیرهای نمادهای غیرمنطقی~$\Lang{L_A}$
\emph{محاسبه‌پذیر}ند. توابع جانشین، جمع و ضرب که به‌ترتیب
$\Assign{\prime}{N}$، $\Assign{+}{N}$ و $\Assign{\times}{N}$اند،
توابعی محاسبه‌پذیر روی اعدادند (برای نمونه، محاسبه‌پذیر با ماشین
تورینگ یا تعریف‌پذیر با بازگشت اولیه). رابطهٔ کوچک‌تر از در
$\Struct{N}$، یعنی $\Assign{<}{N}$، نیز تصمیم‌پذیر است.

مدل‌های نااستاندارد نظریه‌های حسابی چون~$\Th{Q}$ و~$\Th{PA}$ باید
عناصر نااستاندارد داشته باشند. بنابراین دامنه‌هایشان معمولاً افزون
بر~$\Nat$ شامل !!{element}s دیگری نیز هست. بااین‌حال، هر
!!{structure} شمارا را می‌توان روی هر مجموعهٔ !!{denumerable}، از جمله
$\Nat$، بنا کرد. پس مدل‌های نااستانداردی با دامنهٔ~$\Nat$ نیز وجود
دارند. البته در چنین مدل‌هایی چون~$\Struct{M}$، دست‌کم برخی اعداد
نمی‌توانند نقش معمول خود را بازی کنند، زیرا باید عددی چون~$k$ وجود
داشته باشد که برای همهٔ $n \in \Nat$ با
$\Value{\num{n}}{M}$ متفاوت باشد.
\end{explain}

\begin{defn}
!!^a{structure}~$\Struct{M}$ برای~$\Lang{L_A}$
\emph{محاسبه‌پذیر} است هرگاه و تنها هرگاه
$\Domain{M} = \Nat$، توابع $\Assign{\prime}{M}$،
$\Assign{+}{M}$ و $\Assign{\times}{M}$ محاسبه‌پذیر، و رابطهٔ
$\Assign{<}{M}$ تصمیم‌پذیر باشند.
\end{defn}

\begin{ex}\ollabel{ex:comp-model-q}
ساختار~$\Struct{K}$ از
\olref[mdq]{ex:model-K-of-Q} را به یاد آورید. دامنه و تعبیرهای آن چنین
بودند:
$\Domain{K} = \Nat \cup \{a\}$ و
\begin{align*}
  \Assign{\Obj{0}}{K} & = 0\\
  \Assign{\prime}{K}(x) & =
  \begin{cases}
    x+1 & \text{اگر $x\in \Nat$}\\
    a & \text{اگر $x = a$}
  \end{cases}\\
  \Assign{+}{K}(x, y) & =
  \begin{cases}
    x+y & \text{اگر $x$، $y \in\Nat$}\\
    a & \text{در غیر این صورت}
  \end{cases}\\
  \Assign{\times}{K}(x, y) & =
  \begin{cases}
    xy & \text{اگر $x$، $y \in\Nat$}\\
    0 & \text{اگر $x=0$ یا $y=0$}\\
    a & \text{در غیر این صورت}\\
  \end{cases}\\
  \Assign{<}{K} & =
  \Setabs{\tuple{x,y}}{x, y \in \Nat \text{ و } x<y} \cup
  \Setabs{\tuple{x,a}}{x \in \Domain{K}}.
\end{align*}
اما $\Domain{K}$ !!{denumerable} است و بنابراین با~$\Nat$ هم‌شمار
است. برای نمونه، تابع $g\colon \Nat \to \Domain{K}$ که با
$g(0) = a$ و $g(n) = n-1$ برای $n>0$ تعریف می‌شود، !!a{bijection}
است. می‌توانیم آن را به همریختی میان مدلی تازه چون~$\Struct{K'}$
از~$\Th{Q}$ و~$\Struct{K}$ بدل کنیم. در~$\Struct{K'}$ باید توابع و
روابط متفاوتی به نمادهای~$\Lang{L_A}$ نسبت دهیم، زیرا !!{element}s
متفاوتی از~$\Nat$ نقش اعداد استاندارد و نااستاندارد را بازی می‌کنند.

به‌طور مشخص، اکنون~$0$ نقش~$a$ را بازی می‌کند، نه نقش کوچک‌ترین عدد
استاندارد. کوچک‌ترین عدد استاندارد اکنون~$1$ است. پس
$\Assign{\Obj{0}}{K'} = 1$ قرار می‌دهیم. تابع جانشین نیز اکنون متفاوت
است: وقتی عددی استاندارد، یعنی $n>0$، می‌گیرد، هنوز $n+1$ را
بازمی‌گرداند. اما اکنون~$0$ نقش~$a$ را دارد که جانشین خودش است. پس
$\Assign{\prime}{K'}(0) = 0$. برای جمع و ضرب نیز به همین ترتیب داریم
\begin{align*}
\Assign{+}{K'}(x, y) & =
  \begin{cases}
    x+y-1 & \text{اگر $x$، $y >0$}\\
    0 & \text{در غیر این صورت}
  \end{cases}\\
  \Assign{\times}{K'}(x, y) & =
  \begin{cases}
    1 & \text{اگر $x = 1$ یا $y = 1$}\\
    xy - x - y + 2 & \text{اگر $x$، $y > 1$}\\
    0 & \text{در غیر این صورت}\\
  \end{cases}.
\end{align*}
همچنین $\tuple{x, y} \in \Assign{<}{K'}$ هرگاه و تنها هرگاه
$x < y$ و $x > 0$ و $y > 0$، یا $y = 0$.

همهٔ این توابع، توابعی محاسبه‌پذیر روی اعداد طبیعی‌اند و
$\Assign{<}{K'}$ رابطه‌ای تصمیم‌پذیر روی~$\Nat$ است؛ اما آن‌ها همان
توابع جانشین، جمع و ضرب روی~$\Nat$ نیستند و $\Assign{<}{K'}$ نیز همان
رابطهٔ~$<$ روی~$\Nat$ نیست.
\end{ex}

\begin{prob}
!!a{structure}~$\Struct{L'}$ با $\Domain{L'} = \Nat$ ارائه کنید که
با~$\Struct{L}$ از
\olref[mod][mar][mdq]{ex:model-L-of-Q} همریخت باشد.
\end{prob}

\begin{explain}
\Olref{ex:comp-model-q} نشان می‌دهد $\Th{Q}$ مدل‌های محاسبه‌پذیر و
نااستانداردی با دامنهٔ~$\Nat$ دارد. اما نتیجهٔ زیر نشان می‌دهد این امر
برای مدل‌های~$\Th{PA}$ (و ازاین‌رو برای مدل‌های~$\Th{TA}$) صادق نیست.
\end{explain}

\begin{thm}[قضیهٔ تننباوم]
هر مدل محاسبه‌پذیر~$\Th{PA}$ استاندارد است؛ به بیان دیگر، تا حد
همریختی، $\Struct{N}$ تنها مدل محاسبه‌پذیر~$\Th{PA}$ است.
\end{thm}

\end{document}
